\begin{flushleft}
    {\Huge\bfseries Abstract}
\end{flushleft}

\addcontentsline{toc}{chapter}{Abstract}
\vspace{1.5cm}

Emotion detection from social media text is a complex task exacerbated by informal language,
class imbalance, and the inherent ambiguity of human expression. Traditional classification
systems typically enforce a definitive prediction even when the input is semantically unclear
or resides at the boundary between emotional categories. This project introduces a novel
emotion detection framework that incorporates a dedicated ambiguity detection mechanism.
We leverage Topological Data Analysis (TDA), specifically Vietoris-Rips persistent homology,
to analyze the local geometric structure of sentence embeddings surrounding each query. By
integrating this topological complexity—quantified through H1 loops—with the classifier’s
confidence gap, the system can reliably assign an “Ambiguous” label when appropriate, thereby
abstaining from misleading predictions. The system utilizes a hybrid approach combining TF-
IDF features alongside all-mpnet-base-v2 embeddings, classified via one-vs-rest XGBoost
models. Evaluated on the GoEmotions dataset, our approach effectively captures both explicit
emotional signals and genuine semantic ambiguity without resorting to synthetic oversampling,
offering a robust and interpretable framework for natural language understanding.
